\section{Additional steps in the opening twenty-two scans of cote 142}
\label{sec:Z142}

\begin{result}{Z142a}{Proved}{Orientation signs and compactness of the Hom spaces}
For the full, possibly orientation-reversing, finite-cover isotopy
category, the orbits under $M^+(X)$ are indexed by
$R_X(Y)\times\{+1,-1\}$, where $R_X(Y)$ is the finite set of
unoriented isomorphism classes of covers of $Y$ homeomorphic to $X$.
In the oriented completed category, each Hom space is compact and
metrizable. The unqualified assertion that it has no isolated points
is false: for $X=Y=\mathbf P^1(\mathbf C)\setminus\{0,1,\infty\}$
the Hom space consists of the six elements of $S_3$. With individually
labelled punctures it consists of one element. More generally, a
nonempty completed right orbit has no isolated points exactly when
$\widehat{M^+(X)}$ is infinite.
\passage{Cote 142, manuscript pp.~3--6, scans 8--9 and 13.}
\end{result}
\begin{proof}
Every connected cover of an oriented surface has the pulled-back
orientation. A homeomorphism from $X$ to its source has one of two
constant signs. Precomposition by an orientation-preserving map
preserves the sign; two markings of the same sign differ by such a
map. An orientation-reversing homeomorphism of a finite-type oriented
surface exists: realize the compact surface as a symmetric connected
sum of tori, place the punctures on its reflection locus, and transfer
this reflection by a homeomorphism. Both signs therefore occur.
G4 proves that $R_X(Y)$ is finite. G6 and G8 express each oriented
Hom space as a finite union of compact quotients
$S_u\backslash\widehat{M^+(X)}$ with $S_u$ finite.

The mapping-class group is finitely generated: for genus zero,
\cite[Theorem 4.9]{G_FarbMargalit} gives finitely many twists generating
the pure subgroup; adding lifts of generators of the finite puncture
permutation group gives the full group. For positive genus,
\cite[Theorem 4.11]{G_FarbMargalit} gives finite generation of the pure
mapping-class group of $S_{g,n}$ for every $g\geq1,n\geq0$; adding the
finite permutation quotient again gives the full group. A group generated by $r$
elements has only finitely many homomorphisms to each finite group,
since the images of the $r$ generators determine a homomorphism.
There are countably many finite groups. Its profinite completion
therefore has a countable neighbourhood basis and is metrizable,
by the metric determined by any nested cofinal list of open normal
subgroups. Finite unions and quotients by finite groups preserve
compact metrizability.

A point of $S_u\backslash\widehat M$ is isolated precisely when
its finite inverse image is open in $\widehat M$. A profinite group
with a finite open subset has an open identity subgroup and is compact
discrete, hence finite. Conversely a finite quotient space is discrete.
For the three-punctured sphere, cutting along three arcs shows that a
homeomorphism fixing all three punctures is isotopic to the identity:
the complementary discs have fixed boundary after straightening the
arcs, and the disc isotopy is $h_s(z)=s h(z/s)$ on $|z|\le s$, with
identity outside, in a disc coordinate. Thus its full oriented mapping
class group is $S_3$. Euler characteristic forces every self-cover to
have degree one. Its completed Hom space is consequently $S_3$.
The finite permutation calculation is recorded in \texttt{code/G\_checks.py}.
\end{proof}

\subsection{An algebraic description in characteristic zero}
We spell out the path construction proposed in scans 13--16, including
the correspondence that transports a finite cover. For a hyperbolic
type $a=(g,n)$, let $\mathcal M_a$ be the Deligne--Mumford stack of
smooth proper genus-$g$ curves with a reduced unordered divisor of
$n$ points. It is the quotient of the ordered marked-curve stack by
$S_n$. A geometric point is the same datum as a smooth open curve of
type $a$. We use the \emph{\'etale fundamental groupoid}: arrows between
geometric points are natural isomorphisms of their fibre functors on
representable finite \'etale covers. A path has a profinite topology,
being a torsor under either endpoint's profinite fundamental group.

For fixed degree and source and target types, let $\mathcal H$ be the
stack of connected finite \'etale maps of the open curves. Equivalently,
its objects are finite maps of the proper curves whose ramification is
contained in the target marked divisor, with the source marked divisor
its full inverse image. It has maps
\[
       s:\mathcal H\longrightarrow\mathcal M_a,
       \qquad t:\mathcal H\longrightarrow\mathcal M_b .       \tag{Z.1}
\]
The arrows in this stack retain isomorphisms of covers; a cover with
a deck group consequently retains that stabilizer. This convention
specifies the groupoid in (Z.1), including its finite automorphisms.

\begin{result}{Z142b}{Known}{The characteristic-zero algebraic category and its comparison}
Over an algebraically closed field $k$ of characteristic zero there is
a category $\mathcal T_k$ described by the following arrows:
\[
 X\longrightarrow Y:\quad
 (u:X'\longrightarrow Y,\ \alpha:X\rightsquigarrow X'
                         \text{ in }\Pi_1^{\acute et}(\mathcal M_a)),
                                                               \tag{Z.2}
\]
modulo isomorphisms of finite covers that carry the marking $\alpha$
to the other marking. Composition transports the first cover along
the second marking using (Z.1), and then composes the finite maps.
Over $\mathbf C$ this category is naturally equivalent, with its
profinite Hom topologies, to the category of G8 restricted to complex
algebraic curves. Base extension along any embedding $k\hookrightarrow
k'$ of algebraically closed characteristic-zero fields induces an
equivalence $\mathcal T_k\to\mathcal T_{k'}$. Its object isomorphism
classes are precisely the hyperbolic types $(g,n)$, and its
automorphism group at a complex object is the full profinite
completion of the oriented mapping-class group.
\passage{Cote 142, manuscript pp.~6--7, scans 13--16.}
\end{result}
\begin{proof}
We specify and verify the literature inputs. In the opened January
2007 version of Bertin--Romagny \cite{G_BertinRomagny2007}, Theorem 6.9
constructs the marked tame Hurwitz stacks as smooth proper
Deligne--Mumford stacks over $\mathbf Z[1/|G|]$; their smooth-curve
loci are the loci used here. Proposition 6.14 and its proof identify
the target discriminant map as \'etale over smooth pointed curves.
Theorem 6.20 constructs the nongalois-cover stack by quotienting the
Galois-cover stack of fixed monodromy data $(G,H)$, with $H$ core-free,
by the finite group of the monodromy data. These hypotheses apply:
we are in characteristic zero, each Galois closure has finite group
$G$, the subgroup fixing the connected cover is core-free, and all
branch points lie in the marked target divisor. Additional unramified
marked points are allowed by Definition 6.8. Hyperbolicity makes both
marked curves stable. Thus (Z.1) exists algebraically and is \'etale
locally over the target a finite groupoid of covers with finite deck
stabilizers. Finite degree gives finitely many such covers over a
complex target, by G4. The quotient description retains deck groups;
we do not replace this groupoid by its set of isomorphism classes.

Over $\mathbf C$, take a marked oriented surface $S$ of type $b$.
Its Teichm\"uller space is contractible: Fenchel--Nielsen coordinates
identify it with $\mathbf R^{6g-6+2n}$ for every hyperbolic type
\cite[Theorem 10.6 and section 10.6.3]{G_FarbMargalit}. A marking is
an orientation-preserving homeomorphism modulo isotopy. Changing that
marking is exactly the $M(S)$-action, so the associated analytic
moduli stack is $[\operatorname{Teich}(S)/M(S)]$.
Here the quotient-stack assertion can be checked on an arbitrary
family: local topological trivializations give markings, two such
trivializations differ by a mapping class, and their change maps
satisfy the descent cocycle. Conversely the universal marked family
and that descent cocycle recover the family. The contractible space
therefore gives the topological fundamental group $M(S)$.

Fix a topological cover $u:X\to Y$. Pulling a complex structure on
$Y$ through $u$ gives the unique complex structure on $X$ making $u$
holomorphic, including the local completion $z\mapsto z^e$ at every
end. Changing the two markings preserves this cover precisely for
\[
        G_u=\{(a,b)\in M(X)\times M(Y):bu=ua\}.
\]
The corresponding analytic Hurwitz component is consequently
$[\operatorname{Teich}(Y)/G_u]$, with the action through the target
projection. Its fundamental group is $G_u$, and its maps on fundamental
groups are exactly the two projections in G8. Noohi's opened
\cite[Theorem 20.4]{G_Noohi2005} states that, for an algebraic stack
locally of finite type over $\mathbf C$, representable finite \'etale
covers correspond to finite covering stacks of the associated
topological stack. Hence its \'etale fundamental group is the profinite
completion of this topological fundamental group. Applied to the two
stacks just described, this identifies their groups and maps with
$\widehat{M(S)}$ and $\widehat G_u$ and its endpoint projections.
G8 proves the injectivity of the latter stabilizer in the product;
thus the finite deck stabilizer is retained under completion.

SGA~1 XII, Theorem 5.1 algebraizes each finite topological cover of
$Y$. By the preceding path identification, (Z.2) has Hom space
$\coprod_u S_u\backslash\widehat{M(X)}$, exactly (G.1).
For explicit composition, write the second arrow as $(v:Y'\to Z,
\beta:Y\rightsquigarrow Y')$. Lift $\beta$ in the finite-cover
groupoid (Z.1), starting at $u$. It ends at $u':X''\to Y'$ and
produces a source path $\alpha':X'\rightsquigarrow X''$. The product
is $(v u',\alpha'\alpha)$. Existence of the lift follows either from
the locally finite \'etale groupoid description or, over $\mathbf C$,
from the open target image $\widehat H_u$ of $\widehat G_u$ and
its finite cosets. Two lifts differ by the unique deck comparison
respecting their markings, by G8; they give the same arrow. This is
precisely (G.3). Its associativity, continuity, and identity laws were
proved there. The identification proves the claimed comparison.

For base extension we use the opened Bresciani
\cite[Lemma A.22 and Proposition A.23]{G_Bresciani2021}. They give
invariance of the \'etale fundamental gerbe of a geometrically connected
concentrated fibred category under characteristic-zero field extension;
the lemma identifies maps to every finite \'etale stack over an
algebraically closed extension. Our moduli and fixed-degree Hurwitz
stacks are of finite type with quasi-compact separated atlases and
therefore satisfy the concentration hypothesis. We apply the
geometrically connected statement to each component of a fixed-degree
Hurwitz stack separately. The lemma identifies
their fundamental groupoids and all path torsors used in (Z.2).
Applied to the curve $Y$, it also identifies maps to every finite
classifying stack $BG$. Every connected finite cover is a quotient
of its Galois closure; hence the category of finite covers of $Y$
is unchanged. In particular the covers $u$ in (Z.2) descend, together
with their isomorphisms. These identifications preserve composition because the maps commute with
base change. The resulting functor is fully faithful. Every point
over $k'$ belongs to a type component containing a point over $k$;
that component is connected, as its complex Teichm\"uller description
shows after descent to a finitely generated characteristic-zero field.
The path torsor between those points is nonempty: its finite quotient
torsors have the finite intersection property, and their inverse
limit is compact. It supplies an isomorphism in (Z.2), proving essential
surjectivity. This argument identifies types, rather than asserting
that the individual algebraic curves descend up to scheme isomorphism.
Conservativity from G6 ensures that no arrow of degree greater than
one becomes an isomorphism. Degree-one markings give all isomorphisms
within a type. This proves every remaining assertion.
\end{proof}

\begin{result}{Z142c}{Refuted}{The positive-characteristic base-invariance premise}
The assertion on scan 16 that the full geometric fundamental group
of a curve is invariant under every extension of algebraically closed
fields is false in characteristic $p>0$. The counterexample already
occurs on the hyperbolic curve
$U=\mathbf P^1\setminus\{0,1,\infty\}$.
The valid replacements are: invariance for curves in characteristic
zero; invariance for proper curves in arbitrary characteristic; and
invariance of the maximal prime-to-$p$ geometric fundamental group
of a smooth curve in characteristic $p$.
\passage{Cote 142, manuscript p.~7, scans 15--16, parenthetical justification.}
\end{result}
\begin{proof}
Let $k=\overline{\mathbf F}_p$, let $t$ be transcendental over $k$,
and let $k'$ be an algebraic closure of $k(t)$. The equation
\[
                      y^p-y=t x                                      \tag{Z.3}
\]
defines a degree-$p$ finite \'etale $\mathbf Z/p$-torsor over $U_{k'}$:
the polynomial is monic in $y$ and its derivative is $-1$.
Its class does not lie in the image of
$H^1_{\acute et}(U,\mathbf Z/p)\to
H^1_{\acute et}(U_{k'},\mathbf Z/p)$.
Indeed Artin--Schreier theory expresses a torsor by a rational function,
modulo $h^p-h$. A rational function can be reduced at every pole by
successively removing terms whose pole orders are divisible by $p$,
using the $p$th roots of their leading coefficients. Over $k$ all those
roots remain in $k$. The reduced nonconstant principal parts have
orders prime to $p$, and they are unique: a nonconstant $h$ has a pole
of order $m>0$, so $h^p-h$ has a pole of order $pm$, whereas the largest
pole order in a nonzero reduced principal part is prime to $p$.
Constants disappear over an algebraically closed field.
The unique reduced principal part of (Z.3) at infinity is $t x$,
whose coefficient lies outside $k$. Thus it cannot come from a class
over $k$. It is nonzero by the same pole-order argument, so the torsor
is connected. Even changing its chosen $\mathbf Z/p$-action multiplies
$t$ by an element of $\mathbf F_p^*$ and does not make it lie in $k$.

The natural continuous map
$\pi_1(U_{k'})\to\pi_1(U_k)$ would induce a bijection on continuous
maps to $\mathbf Z/p$ if it were an isomorphism. The displayed
cohomology failure proves that it is not an isomorphism.
For the replacements, SGA~1 X, Corollary 1.8 gives invariance for
proper schemes over algebraically closed fields, and XIII,
Corollaries 2.9 and 2.12 give the prime-to-characteristic presentation
and comparison for smooth curves \cite{SGA1}. In characteristic zero
the same presentation is the full one; equivalently the finite-cover
comparison and descent used in Z142b apply. These statements have
exactly the indicated properness or prime-to-$p$ hypotheses.
The counterexample refutes the stated premise for the proposed wild
category. It does not identify an unspecified replacement for that
category: its required path-transport construction is recorded
separately in Z142g below.
\end{proof}

\begin{result}{Z142d}{Proved}{Field automorphisms and the actual outer ambiguity}
Every automorphism $\sigma$ of an algebraically closed characteristic-zero
field $k$ induces a functor on $\mathcal T_k$ by base change, and these
functors compose with their canonical base-change isomorphisms. After
choosing one object of each type and paths to its conjugates, the
induced automorphisms of its profinite automorphism group are defined
up to inner automorphisms. Changing a path by $a$ changes the induced
automorphism by conjugation by $a$.
\passage{Cote 142, manuscript pp.~7--8, scans 16--17.}
\end{result}
\begin{proof}
Send $X$ to $X^\sigma$, the cover $u$ to $u^\sigma$, and a natural
isomorphism of fibre functors to its pullback along $\sigma$.
Base change preserves finite \'etale covers and the stack maps (Z.1);
it therefore preserves the composition already constructed. The
canonical isomorphism $(X^\sigma)^\tau\simeq X^{\tau\sigma}$, with
the order fixed by this formula, satisfies the cocycle identity by
the associativity of tensor products. These maps give the promised
action without identifying a curve with its conjugate.
A chosen path $c:X\rightsquigarrow X^\sigma$ identifies the endpoint
groups by $b\mapsto c^{-1}b c$. Replacing $c$ by $ca$ conjugates this
identification by $a^{-1}$. This explicit formula proves the claimed
ambiguity and gives a well-defined outer action.
\end{proof}

\begin{result}{Z142e}{Refuted}{Path isotopy after allowing extra omitted branch points}
For maps obtained by composing finite covers and inclusions that fill
punctures, a continuous path through \'etale maps need not lie in a
single two-sided orbit of the identity components of the source and
target homeomorphism groups. The following example is an explicit
composite of the two allowed sorts of arrows, has degree six,
and has fixed source type $(0,20)$ and target type $(0,3)$.
For a continuous family of holomorphic maps of compact curves with
fixed ramification data and fixed incidence partition of the images
of all marked points, local path isotopies do lie in such a two-sided
orbit. The incidence hypothesis is essential in the example.
\passage{Cote 142, manuscript pp.~8--9, scans 18--19; ``isotopie raisonnable''.}
\end{result}
\begin{proof}
For $0\leq t\leq1/10$ set
\[
 h_t(z)=z^3-3z+t,\qquad f_t(z)=h_t(z)^2,
 \qquad Y=\mathbf P^1\setminus\{9,16,\infty\}.
\]
Define the intermediate curves and the source by
\[
 \begin{aligned}
 Y_3&=\mathbf P^1\setminus\{0,9,16,\infty\},\\
 Y_2&=\mathbf P^1\setminus\{0,3,-3,4,-4,\infty\},\\
 Y_{1,t}&=Y_2\setminus\{t-2,t+2\},\qquad
 X_t=h_t^{-1}(Y_{1,t}).
 \end{aligned}
\]
The restriction of $f_t$ is the actual composite
\[
 X_t\xrightarrow[\text{finite \'etale, degree }3]{h_t}Y_{1,t}
 \hookrightarrow Y_2
 \xrightarrow[\text{finite \'etale, degree }2]{w\mapsto w^2}Y_3
 \hookrightarrow Y.                                      \tag{Z.4}
\]
All five curves are hyperbolic. The cubic has finite critical points
$1,-1$ and critical values $t-2,t+2$, both removed from $Y_{1,t}$;
infinity is removed as well. This proves that the first finite
map is \'etale. The removed sets for the square map are exactly
the inverse images of $0,9,16,\infty$, proving the same assertion
for the second finite map. Each inclusion fills the stated punctures.

The source has twenty punctures: the fifteen simple preimages under
$h_t$ of $0,3,-3,4,-4$, the four points $1,-2,-1,2$, and infinity.
Indeed,
\[
 h_t(z)-(t-2)=(z-1)^2(z+2),\qquad
 h_t(z)-(t+2)=(z+1)^2(z-2).
\]
For each $c\in\{0,3,-3,4,-4\}$ the program checks
\[
 \operatorname{disc}_z(h_t-c)=27\bigl(4-(t-c)^2\bigr).
\]
These five discriminants never vanish on the chosen interval:
$|t|<2$, whereas $|t-c|>2$ for the four nonzero values of $c$.
Also $t-2,t+2$ are distinct and avoid that five-element set.
Thus all twenty punctures remain distinct. Tracking them and
applying disjoint-disc isotopies gives homeomorphisms
$\eta_t:X_0\to X_t$ with $\eta_0=1$. The maps
$f_t\eta_t:X_0\to Y$ form a continuous path of local
homeomorphisms. They belong to the category generated by the
arrows in (Z.4), including the source homeomorphism $\eta_t$.

The derivative is $f'_t=6h_t(z)(z^2-1)$. Consequently its finite
critical values are exactly
\[
 0,\quad(t-2)^2,\quad(t+2)^2.
\]
The last two differ by $8t$. At $t=0$ there are two distinct
finite critical values, $0,4$; at every $0<t\leq1/10$ there
are three. This calculation and the displayed factorizations
are checked exactly in \texttt{code/G\_checks.py}.
A homeomorphism of a finite-type surface extends uniquely to the
finite-end compactification. A two-sided conjugacy of two maps on
these punctured surfaces therefore extends to their compact spheres
and carries their sets of critical values bijectively. The two maps
cannot be in the same two-sided orbit, because those sets have
different cardinalities (including infinity adds one on both sides).
This is an explicit failure for the ordinary path interpretation of
``isotopie''. The manuscript does not define the adjective
``raisonnable'', so this example does not select a final reading of
that adjective.

For the corrected local assertion, track the finitely many distinct
marked images and branch values in disjoint discs; the fixed incidence
partition ensures that the same labels remain together. An ambient
isotopy of the target carries all these points back to their initial
positions. Away from those values the family is a family of finite
unramified covers of a fixed punctured compact curve. The implicit
function theorem gives the covering property also in the parameter
direction; properness over compact subsets makes it a finite cover
of the product with an interval. Path lifting in the interval
trivializes that cover. Extend this trivialization over each omitted
point using its unchanged local degree $z\mapsto z^e$. The fixed
incidence data tracks the remaining marked preimages as well. It
gives a source isotopy and the desired two-sided equivalence.
Thus fixed Hurwitz data supplies the exact transfer map that an
arbitrary path in the example lacks.
\end{proof}

\begin{result}{Z142f}{Proved}{Filling punctures and the necessary source containment}
For finite maps $\overline f:\overline X\to\overline Y$ of smooth
proper curves in characteristic zero, and finite reduced sets $S,T$,
the restriction $X=\overline X\setminus S\to
Y=\overline Y\setminus T$ is defined and \'etale exactly when
$\overline f^{-1}(T)\subset S$ and every ramification point of
$\overline f$ lies in $S$. Composing finite \'etale maps and puncture
fillings produces such maps, including maps whose omitted
ramification points have images inside $Y$.
For ordered punctures, forgetting a specified set $I$ induces a
functor on \'etale fundamental groupoids of the corresponding moduli
stacks whenever the remaining type is hyperbolic. Successive
forgettings compose with their canonical natural isomorphisms.
\passage{Cote 142, manuscript pp.~8--10, scans 18, 20--21.}
\end{result}
\begin{proof}
The first containment is exactly the condition for the restriction to
land in $Y$. In local parameters, a finite map of smooth curves has
$y=x^e$ times a unit; in characteristic zero it is \'etale at that
point exactly when $e=1$. This proves the second condition and both
necessity and sufficiency. For a concrete composite, take
\[
 \mathbf P^1\setminus\{0,\infty,1,-1,i,-i\}
   \xrightarrow{z\mapsto z^2}
 \mathbf P^1\setminus\{0,1,-1,\infty\}
   \hookrightarrow
 \mathbf P^1\setminus\{1,-1,\infty\}.
\]
The first arrow is finite \'etale of degree two; the second fills the
puncture $0$. Its composite omits the ramification point $0$ although
its branch value $0$ is inside the final target. All three open curves
are hyperbolic. The inverse images and derivative are checked in
\texttt{code/G\_checks.py}.

On families of smooth proper marked curves, the forgetting map keeps
the underlying family and the remaining disjoint sections. With a
hyperbolic remaining type these are again objects of the indicated
moduli stack, and the operation commutes with every base change and
isomorphism of families. This defines an algebraic morphism of stacks.
Pullback of representable finite \'etale covers defines its functor on
fibre functors and their natural isomorphisms, hence on the \'etale
fundamental groupoids. The composite forgetting two disjoint sets is
the same operation on families as forgetting their union; the
resulting natural isomorphisms satisfy associativity. Unordered
remaining points are handled by their finite permutation quotient.
The hyperbolicity condition is necessary to stay in the category
under discussion: forgetting one puncture of type $(0,3)$ produces
type $(0,2)$. This gives the exact type restriction omitted in scan 20.
\end{proof}

\begin{result}{Z142g}{Open}{The uncompleted parts of the later programme}
The following objectives in scans 15--21 are not established here:
(i) a wild positive-characteristic version of (Z.2) with the stated
finite-cover factorization, path transport, and base-extension
equivalence; (ii) the requested intrinsic description of the entire
category together with the image of the arithmetic action, using
only the specified profinite and combinatorial data; (iii) a
completion theorem for the enlarged category containing arbitrary
puncture fillings, with the unspecified ``reasonable'' isotopy of
scan 19.
\passage{Cote 142, manuscript pp.~7--10, scans 15--21.}
\end{result}
\begin{proof}[Recorded attacks and exact obstructions]
For (i), we constructed the characteristic-zero category and its
comparison in Z142b, then tested the asserted base-invariance premise
on degree-$p$ covers of $U_{0,3}$. Equation (Z.3) proves that this
premise fails. The next required datum is a replacement path-lifting
rule for the wild Hurwitz stack, compatible with its families of
covers of a fixed target; the finite discrete transport argument in
Z142b does not give that rule. We have not identified a category
satisfying all three requested properties in this setting.

For (ii), Z142d constructs the arithmetic functors and explicitly
computes their inner ambiguity, and G6--G8 give the entire completed
finite-cover category from the discrete surface category. These
constructions supply no criterion deciding which compatible
profinite automorphisms are induced by field automorphisms. The
remaining task is an image computation for the displayed base-change
action. Completing mapping-class groups does not compute that image:
the functor $\sigma\mapsto(-)^\sigma$ is additional data, and its
formula in Z142d still uses the field. No surjectivity claim about
this action is made in this paper.

For (iii), we first tested ordinary path isotopy; Z142e refutes its
identification with two-sided isotopy. We then proved the local
comparison with fixed Hurwitz incidence data, and Z142f constructs
the actual puncture-forgetting functors. These do not specify the
relations among arbitrary composites when branch-image incidence
changes. A free category on these functors can be constructed as the
set of finite composable words, modulo chosen relations, but no
relations were given in the manuscript that imply the proposed
comparison. The next mathematical step is an explicit choice of the
isotopy relation and a theorem identifying its groupoid on each
ramification stratum and across its boundary. We retain this missing
comparison as an exact unfinished objective; we do not turn the
programmatic wording into an asserted theorem.
\end{proof}
